\printedpage{124}
% head: A. Grothendieck : Dérivateurs, Ch. XVIII, Édition des Universités de Montpellier II et Paris VII

compatibilities $3_\alpha(\beta)$, hence double arrows
\[
\begin{tikzcd}
\varphi(i_\beta) \arrow[r, shift left, "\tau'_{\alpha,\beta}u_\beta"] \arrow[r, shift right, "u_\alpha\varphi(\tau_{\alpha,\beta})"'] & j_\alpha \; .
\end{tikzcd}
\]
As the set of these double arrows has cardinality $\card\alpha < \pi$, and $J'$ is large relative to the said cardinal, we
\msp{page 329}
thus find an arrow
\[
\rho : j_\alpha \longrightarrow j
\]
in $J$ which equalizes them all. As $J'$ is cofinal, we can suppose $j$ in $J'$, say $\bar{j}_\alpha$. Thus, at the cost of replacing $j_\alpha$ by $\bar{j}_\alpha$, and the $\tau'_{\alpha,\beta}$ by $\bar{\tau}'_{\alpha,\beta} = \rho\tau'_{\alpha,\beta}$, finally $u_\alpha : \varphi(i_\alpha) \longrightarrow j_\alpha$ by $\bar{u}_\alpha = \rho u_\alpha : \varphi(i_\alpha) \longrightarrow \bar{j}_\alpha$, we satisfy all the desired constraints.

I summarize the result (``$\pm$ tautological''!) proved on pages 313-329.

% a vertical change bar is printed in the left margin alongside the statement from "Lemme-clef" to "suivantes :"
\textbf{Key lemma 4.13.17.} \textit{Let us consider a cartesian square in} Cat
\[
\begin{tikzcd}
I \arrow[r, "\varphi"] & J \\
I' \arrow[u, hook] \arrow[r] & J' \arrow[u, hook, "\psi"'] \; .
\end{tikzcd}
\]

\textit{We make the following assumptions:}

\msp{page 330}

\begin{enumerate}[label={}]
\item[\textit{1)}] \textit{$\varphi$ and $\psi$ are cofinal, $\psi$ is the inclusion of a strictly full subcategory $J'$ of $J$, $I$ is filtering (hence $J$ and $J'$ as well).}
\end{enumerate}

\textit{We seek conditions for the inclusion of the strictly full subcategory $I' \longrightarrow I$ to be cofinal (in which case we see at once that $I'$ is filtering, and $I' \longrightarrow J'$ is also cofinal). For this, we suppose that there exists a \textup{limit ordinal} $\alpha_0$ having the following additional properties:}

\begin{enumerate}[label={}]
\item[\textit{2)}] \textit{$I$ and $J'$ are stable under $\varinjlim$ of type $I_{\alpha_0}$ $(= \{\beta \mid \beta \ \textit{ordinal}\ < \alpha_0\})$, and the functors $\varphi$, $\psi$ commute with these limits.}
\item[\textit{3)}] \textit{For every $\alpha \in I_{\alpha_0}$, $I$ is large relative to $\card\alpha$. (Thus, $\varphi$ being cofinal, $J$ is also large relative to $\card\alpha$, and likewise, $\psi$ being fully faithful, and cofinal, $J'$ is large relative to $\card\alpha$.)}
\end{enumerate}

\textit{Under these conditions, $I'$ is cofinal in $I$ (hence also large relative to $\card\alpha$, for every $\alpha \in I_{\alpha_0}$).}

The proof is the one given above. We did not have to use that the ordinal $\alpha_0 = \alpha_\pi$ was of this particular kind, only that it was a limit cardinal.

\printedpage{125}
% head: A. Grothendieck : Dérivateurs, Ch. XVIII, Édition des Universités de Montpellier II et Paris VII

\msp{page 331}

This lemma implies, as we saw on page 312, the implication b $\Longrightarrow$ d in the diagram of implications $(*)$ on page 310, which is thus now proved.

To finish proving theorem 4.13.15, it remains to prove the implications
\[
\begin{tikzcd}[row sep=small]
\mathrm{c}' \arrow[dr, Rightarrow] & & \\
& \mathrm{a}' \arrow[r, Rightarrow] & \mathrm{a} \; . \\
\mathrm{b} \arrow[ur, Rightarrow] & &
\end{tikzcd}
\]

Let us therefore prove (b or c$'$) $\Longrightarrow$ a$'$. It remains to prove that $\underline{S}$ is $\pi$-accessible in $\mathcal{M}$, i.e. that the fully faithful functor
\[
\Ind_\pi(\underline{S}(\pi)) \longrightarrow \mathcal{M} \tag*{$(*)$}
\]
has essential image $\underline{S}$. Thus we must prove that if $x \in \Ob\mathcal{M}$, then
\[
x \in \Ob\underline{S} \quad \overset{(?)}{\Longleftrightarrow} \quad
\begin{minipage}[c]{0.55\textwidth}
$x$ is in the essential image of $\Ind_\pi(\underline{S}(\pi))$, i.e. $x$ is isomorphic to $\varinjlim\limits_{i\in I} x_i$, the $x_i$ in $\Ob\underline{S}(\pi)$, $I$ large relative to $\pi$.
\end{minipage}
\tag*{$(*)$}
\]

\msp{page 332}

$\Bigg($ Let us first suppose that $\underline{S}$ satisfies b) (\footnote[141]{Clumsy, we can treat the two cases b) and c) at the same time, cf. bottom of page 333.}). Thus we know that
\[
x \in \Ob\underline{S} \qquad \Longleftrightarrow \qquad \underline{S}/x \text{ is cofinal in } \mathcal{M}_\pi/x \; . \tag*{$(**)$}
\]

We want to deduce the equivalence $(*)$ from the equivalence $(**)$.

Let us prove $\Longrightarrow$ in $(*)$, i.e. let us suppose $x \in \Ob\underline{S}$, and prove that $x$ is in the essential image of $\Ind_\pi\underline{S}(\pi)$. By virtue of $(**)$, $\underline{S}(\pi)/x$ is a subcategory (strictly full) cofinal in $\mathcal{M}_\pi/x$, thus since $\mathcal{M}_\pi/x$ is large relative to $\pi$, $\underline{S}(\pi)/x$ is as well. As we know that
\[
x \ = \ \varinjlim_{i\in\mathcal{M}_\pi/x} x_i
\]
($\pi$ being a useful cardinal for $\mathcal{M}$), it follows by cofinality that we also have
\[
x \ = \ \varinjlim_{i\in\underline{S}(\pi)/x} x_i \; ,
\]
and since $\underline{S}(\pi)/x$ is large relative to $\pi$, this implies that $x$ belongs to the essential image of $\Ind_\pi(\underline{S}(\pi))$.

Let us prove $\Longleftarrow$ in $(*)$, i.e. let us suppose that $x$ is in the essential image of $\Ind_\pi(\underline{S}(\pi))$,
\printedpage{126}
% head: A. Grothendieck : Dérivateurs, Ch. XVIII, Édition des Universités de Montpellier II et Paris VII
\msp{page 333}
\[
x \ = \ \varinjlim_{i\in I} x_i \; , \qquad\quad \text{the } x_i \text{ in } \underline{S}(\pi),\ I \text{ large relative to } \pi.
\]

Then the composite functor $i \longmapsto x_i$,
\[
I \ \longrightarrow \ \underline{S}(\pi)/x \ \hookrightarrow \ \mathcal{M}_\pi/x
\]
is cofinal (cf. 4.4.2, page 79), hence $\underline{S}(\pi)/x \longrightarrow \mathcal{M}_\pi/x$ satisfies condition F 1 (since the composite satisfies it), and since $\underline{S}(\pi)/x \longrightarrow \mathcal{M}_\pi/x$ is fully faithful, this already implies that this functor is cofinal. Thus by $(**)$ we conclude from this $x \in \Ob\underline{S}$, q.e.d.

Let us now suppose that $S$ satisfies c$'$ (instead of b), and let us again prove the equivalence $(*)$. For this, let us write
\[
\square \qquad x \ = \ \varinjlim_{i\in I} x_i \qquad\quad
\left\{
\begin{minipage}[c]{0.5\textwidth}
the $x_i \in \Ob\mathcal{M}_\pi$, $I$ \textit{ordered set}, $(x_i)_{i\in I}$ is $\pi_0$-$\pi$-adapted (\footnotemark[142]).
\end{minipage}
\right.
\]
\footnotetext[142]{\textbf{N.B.} That exists, thanks to 4.13.14.}

which is possible since $\pi$ is useful for $\mathcal{M}$ $\Bigg)$. Let us define $I_{\underline{S}}$ as in the statement of c), hence as the inverse image of $\underline{S}(\pi)/x \hookrightarrow \mathcal{M}_\pi/x$ under $I \longrightarrow \mathcal{M}_\pi$. By hypothesis d$'$) \texttt{[rather c$'$)]}, we have
\[
x \in \Ob\underline{S} \qquad \Longleftrightarrow \qquad I_{\underline{S}} \text{ is cofinal in } I \; . \tag*{$(**')$}
\]

Using $(**')$, let us again prove $(*)$ (page 331). In fact, we shall use only that we have
\msp{page 334}
the isomorphism $\square$, with $I$ large relative to $\pi$, and that we have the equivalence $(**')$, without having to use that $I$ is ordered. Thus the proof also holds in case c).

Let us prove $\Longrightarrow$ in $(*)$, i.e. let us suppose $x \in \Ob\underline{S}$, let us prove that $x$ is in the essential image of $\Ind_\pi(\underline{S}(\pi))$. By virtue of $(**')$, $I_{\underline{S}}$ is cofinal in $I$, and since $I_{\underline{S}} \longrightarrow I$ is a fully faithful inclusion, it follows that $I_{\underline{S}}$ is also large relative to $\pi$, on the other hand we shall have by cofinality
\[
x \ = \ \varinjlim_{i\in I_{\underline{S}}} x_i \; ,
\]
where this time the $x_i$ are in $\underline{S}(\pi) = \underline{S}\cap\mathcal{M}_\pi$, whence it follows that $x$ is in the essential image of $\Ind_\pi(\underline{S}(\pi))$.

Let us prove $\Longleftarrow$ in $(*)$, i.e. let us suppose that we have
\[
x \ = \ \varinjlim_{j\in J} x_j \; , \qquad\quad \text{the } x_j \text{ in } \underline{S}(\pi),\ J \text{ large relative to } \pi \; .
\]

\printedpage{127}
% head: A. Grothendieck : Dérivateurs, Ch. XVIII, Édition des Universités de Montpellier II et Paris VII

\msp{page 335}

Invoking 4.4.2 once again, we find that $J \longrightarrow \mathcal{M}_\pi/x$ is cofinal, and since this functor factors through the fully faithful inclusion functor $\underline{S}(\pi)/x \hookrightarrow \mathcal{M}_\pi/x$, it follows as before that the latter is cofinal. We shall admit for a moment that this implies that the inverse image $I_{\underline{S}}$ of this strictly full subcategory, under the likewise cofinal functor $I \longrightarrow \mathcal{M}_\pi/x$, is cofinal in $I$. This being so, the equivalence $(**')$ (p. 333) implies $x \in \Ob\underline{S}$, which is what we wanted.

Thus it remains to prove that if
\[
J' \ = \ \underline{S}(\pi)/x \hookrightarrow \mathcal{M}_\pi/x \ = \ J
\]
is cofinal, then $I' = I_{\underline{S}} \longrightarrow I$ is as well. It seems that we really do have to invoke again the stability conditions $\mathrm{L}(\pi_0,\pi)$ made in d$'$) \texttt{[rather c$'$)]} which had not been necessary for case c) and
\msp{page 336}
the (a bit) nasty key lemma 4.13.17, where we shall take (as in the proof written
before the statement of this lemma) % [sic] cet lemme -> ce lemme, demonstrative adjective "cet" before a consonant
$\alpha = \alpha_\pi \overset{\text{def}}{=}$ smallest ordinal $\alpha$ such that $\card\alpha = \pi$. Thus $I_\alpha$ has cardinality $\pi$, and is large relative to every cardinal $\pi' < \pi$, and in particular relative to $\pi'_0$. Assumption $2^\circ$) of 4.13.17 is satisfied for $I$, by definition of a $\pi_0$-$\pi$-adapted direct system, and for $J'$, thanks to assumption $\mathrm{L}(\pi_0,\pi)$ on $\underline{S}(\pi)$. Likewise, assumption $3^\circ$) ($I$ large relative to every cardinal $< \pi$) is satisfied by definition of $\pi_0$-$\pi$-adapted systems -- $I$ is even large relative to $\pi$. Assumption $1^\circ$) has already been seen.

We have thus proved (by a circular route)
\[
\begin{tikzcd}[row sep=small]
& & & \mathrm{c}' & \\
\mathrm{a} \arrow[r, Rightarrow] & \Bigg(\mathrm{a}' \arrow[r, Leftrightarrow] & \mathrm{c} \arrow[ur, Leftrightarrow] \arrow[dr, Rightarrow] & & \Bigg)\, , \\ % Verified against PDF leaf 9: the oblique equivalence has both arrowheads.
& & & \mathrm{b} &
\end{tikzcd}
\]

\texttt{[and also b $\Longrightarrow$ c,]} having to invoke the nasty key lemma for the \textit{two} implications
\[
\mathrm{a}' \ \Longrightarrow \ \mathrm{c} \qquad\quad \text{and} \qquad\quad \mathrm{c}' \ \Longrightarrow \ \mathrm{a}'
\]
(hence for two implications \textit{in opposite directions}, a strange thing \dots).

\msp{page 337}

It only remains still to prove the implication
\[
\mathrm{a}' \ \Longrightarrow \ \mathrm{a} \; ,
\]
which is essentially sorital.

This completes the proof of the accessibility criterion 4.13.15 for a full subcategory $\underline{S}$ of $\mathcal{M}$.

We are going to give a somewhat decanted variant of 4.13.15, which (with slightly weaker assumptions) captures the essential part of its substance. Statement 4.13.15 is an almost immediate consequence of it:

\printedpage{128}
% head: A. Grothendieck : Dérivateurs, Ch. XVIII, Édition des Universités de Montpellier II et Paris VII

\msp{page 338}

\textbf{Corollary 4.13.18} (of 4.13.15). \textit{Let $\mathcal{M}$ be an accessible category, $\pi$ a useful cardinal for $\mathcal{M}$, $\alpha_0$ a \textup{limit ordinal} such that we have $\card\alpha_0 \leqslant \pi$. Let $I_{\alpha_0} = \{\alpha \mid \alpha < \alpha_0\}$, and suppose that $\mathcal{M}$ is stable under $\varinjlim$ of type $I_{\alpha_0}$.}

\textit{Let then $\underline{S}$ be a strictly full subcategory of $\mathcal{M}$, $\underline{S}(\pi) = \underline{S}\cap\mathcal{M}_\pi$. Let us denote by $L(\alpha_0)$ the condition: $\underline{S}(\pi)$ is stable in $\mathcal{M}_\pi$ (or what amounts to the same thing, in $\mathcal{M}$) under $\varinjlim$ of type $I_{\alpha_0}$. Then the following conditions are equivalent.}

\begin{enumerate}[label={}]
\item[\textit{a)}] \textit{$\underline{S}$ is $\pi$-accessible in $\mathcal{M}$, and is stable in $\mathcal{M}$ under $\varinjlim$ of type $I_{\alpha_0}$.}
\item[\textit{a$'$)}] \textit{$\underline{S}$ is $\pi$-accessible and $\underline{S}(\pi)$ satisfies $L(\alpha_0)$.}
\item[\textit{b)}] \textit{$\underline{S}(\pi)$ satisfies $L(\alpha_0)$, and for every $x \in \Ob\mathcal{M}$, we have this: $x \in \Ob\underline{S}$ if and only if the strictly full subcategory $\underline{S}(\pi)/x$ of $\mathcal{M}(\pi)/x$ is cofinal in $\mathcal{M}(\pi)/x$} \texttt{[read $\mathcal{M}_\pi/x$]}.
\item[\textit{c)}] \textit{$\underline{S}(\pi)$ satisfies $L(\alpha_0)$, moreover for every $x \in \Ob\mathcal{M}$, and every representation}
\[
x \ \simeq \ \varinjlim_{i\in I} x_i
\]
\msp{page 339}
\textit{which is ``$\alpha_0$-$\pi$-adapted'', i.e. such that}
\begin{enumerate}[label={}]
\item[\textit{$1^\circ$)}] \textit{$I$ equivalent to a small category, and large relative to $\pi$,}
\item[\textit{$2^\circ$)}] \textit{the $\varinjlim$ of type $I_{\alpha_0}$ exist in $I$,}
\item[\textit{$3^\circ$)}] \textit{the functor $I \longrightarrow \mathcal{M}$, $i \longmapsto x_i$, commutes with the $\varinjlim$ of type $I_{\alpha_0}$, and finally}
\item[\textit{$4^\circ$)}] \textit{the $x_i \in \Ob\mathcal{M}_\pi$}
\end{enumerate}
\textit{(compare 4.13.11), we have this (where we denote by $I_{\underline{S}}$ the strictly full subcategory of $I$ formed by the $i$ such that $x_i \in \Ob\underline{S}$): we have $x \in \Ob\underline{S}$ if and only if $I_{\underline{S}}$ is cofinal in $I$.}
\item[\textit{c$'$)}] \textit{(\textup{If} $I_{\alpha_0}$ is large relative to $\pi_0$, and $\pi^{\pi_0} = \pi$, to be able to invoke 4.13.14, cf. p. 335.) As in c), but assuming $I$ ordered.}
\end{enumerate}
(\footnote[143]{\textbf{N.B.} If we do not suppose $I_{\alpha_0}$ large relative to $\pi_0$ and % [sic] suppose par -> suppose pas, negation "ne ... pas"
$\pi^{\pi_0} = \pi$, then all the conditions from a) to c) (a, a$'$, b, c) are equivalent and \textit{imply} c, but are perhaps not equivalent to it.}).

\textsc{Proof}. It is carried out following the same logical diagram as that of 4.13.15,
\[
\begin{tikzcd}[row sep=small]
& & & \mathrm{c}' \arrow[dr, Rightarrow, "\textbf{!}"] & \nearrow\ \begin{tabular}[b]{@{}r@{}}\footnotesize here, we must suppose\\ \footnotesize $I_{\alpha_0}$ large relative to $\pi_0$\\ \footnotesize and $\pi^{\pi_0} = \pi$.\end{tabular} & \\
\mathrm{a} \arrow[r, Rightarrow] & \mathrm{a}' \arrow[r, Rightarrow, "\textbf{!}"] & \mathrm{c} \arrow[ur, Rightarrow] \arrow[dr, Rightarrow] & & \mathrm{a}' \arrow[r, Rightarrow] & \mathrm{a} \; , \\
& & & \mathrm{b} \arrow[ur, Rightarrow] & &
\end{tikzcd}
\]
where we have marked with the sign\,! the two implications which are ``nasty'', in the sense that they call on the key lemma 4.13.17 (pages 329-330), which is proved by a somewhat twisted transfinite construction. The proof of the seven implications of the diagram is essentially
the one given on the preceding pages 311-337, % [sic] celle donné -> celle donnée, agreement with feminine "celle"